\pdfbookmark[1]{Abstract}{Abstract}
\thispagestyle{empty}
\begin{center}
  \rule{\textwidth}{0.4pt}\\[0.9em]
  {\Large\spacedallcaps{Abstract}}\\[0.7em]
  \rule{\textwidth}{0.4pt}
\end{center}
\bigskip

\noindent Concurrent queues couple synchronization to storage: a compare-and-swap loop or a
per-slot sequence counter guards every slot, and the proof of the queue is the proof of that
coupling. We separate the two. A \emph{turnstile} is a manifold of lanes, each one atomic word on
its own cache line. A reader enters with one fetch-or, leaves with one fetch-add, and a writer
drains by reading each open lane once and waiting on the busy ones. A \emph{ticket ring} is a
power-of-two array of slots that a 64-bit stamp indexes with one mask; its push and pop have no
branch and no atomic operation, because the turnstile is the only guard. Every slot write records
the stamp that wrote it, so the ring is a log of causal events as well as a queue.

The C11 implementation is seven headers. The Lean~4 formalization proves the lane word at 32 and
64 bits, the open count, the queue word, the backoff, the drain and the sleeps of the turnstile, and
the stamp, the place, the window and the pushes of the ring. GenMC checks the drain and the resize
of the manifold under five memory models. We state each theorem as it is checked, we name the
trusted axioms (the atomics of the target, the futex contract of the host, the depth bound that the
grantor of tickets keeps), and we state what the artifact does not provide: a rewind operation, a
parity bit, and a sandbox case study are design arguments over the proved interface, not
machine-checked results.
\cleardoublepage
